\section{Обзор литературы и теоретические основы}

\subsection{Гидродинамическая теория смазки}

Основой анализа подшипников скольжения является гидродинамическая теория смазки, заложенная в работах О.~Рейнольдса~\cite{pinkus}. Несущая способность смазочного слоя возникает вследствие вязкого течения жидкости в клиновом зазоре между вращающимся валом и неподвижной втулкой.

В состоянии покоя вал лежит на дне втулки, и между поверхностями существует непосредственный контакт. При вращении вала смазка за счёт вязкости затягивается в сужающийся зазор между валом и втулкой. В области сужения зазора накапливается давление, которое формирует несущий слой, и вал «всплывает», занимая эксцентричное положение~\cite{hamrock}. Величина эксцентриситета $e$ определяется балансом гидродинамической силы и внешней нагрузки. Относительный эксцентриситет $\varepsilon = e/c$ ($c$ -- радиальный зазор) является ключевым параметром: при $\varepsilon \to 0$ вал центрирован, при $\varepsilon \to 1$ -- приближается к контакту с втулкой~\cite{pinkus}.

В зависимости от толщины смазочного слоя различают три режима смазки, описываемые кривой Штрибека~\cite{hamrock}. Граничная смазка реализуется при малых скоростях или высоких нагрузках, когда толщина плёнки сопоставима с шероховатостью поверхностей и нагрузка передаётся через контакт микронеровностей. Смешанная смазка представляет переходный режим, в котором часть нагрузки воспринимается гидродинамическим давлением, а часть -- контактом шероховатостей. Жидкостная (гидродинамическая) смазка характеризуется полным разделением поверхностей непрерывной плёнкой смазки, и коэффициент трения определяется только вязкостью жидкости~\cite{pinkus, hamrock}.

Величина радиального зазора $c$ между валом и втулкой является важным конструктивным параметром. При слишком малом зазоре затруднён подвод смазки, возрастает риск сухого трения и заедания. При слишком большом зазоре несущая способность падает, так как давление в расширенном зазоре недостаточно для поддержания вала~\cite{chernjavsky}. Оптимальный зазор определяется из условия обеспечения минимального коэффициента трения при заданной нагрузке.

Основные интегральные характеристики подшипника скольжения~\cite{hamrock}:

Несущая способность $F$ -- равнодействующая гидродинамического давления, определяемая как интеграл давления по поверхности подшипника. Нагрузочная способность возрастает с увеличением вязкости смазки, скорости вращения и уменьшением зазора.

Коэффициент трения $\mu$ -- отношение силы вязкого трения к нагрузке. В режиме жидкостной смазки $\mu$ определяется градиентом скорости в смазочном слое и обычно составляет $10^{-3}\text{ - }10^{-2}$~\cite{pinkus}.

Расход смазки $Q$ -- объёмный поток жидкости через торцы подшипника (торцевые утечки). Расход определяет интенсивность обновления смазки в зазоре и, следовательно, условия теплоотвода~\cite{hamrock}.

Основные допущения, принимаемые при выводе уравнения Рейнольдса~\cite{hamrock}: смазка -- несжимаемая ньютоновская жидкость; течение ламинарное; инерционные силы пренебрежимо малы по сравнению с вязкими; толщина смазочного слоя мала по сравнению с радиусом ($c \ll R$); давление постоянно по толщине плёнки; отсутствует проскальзывание жидкости на границах.

\subsection{Дискретный микрорельеф поверхностей трения}

Одним из перспективных методов улучшения трибологических характеристик подшипников скольжения является нанесение на рабочую поверхность дискретного микрорельефа -- системы регулярных углублений микронного масштаба. Углубления формируются на поверхности втулки или вала с заданной геометрией, глубиной и шагом размещения~\cite{etsion2005}. Идея текстурирования основана на том, что каждое микроуглубление способно генерировать дополнительное гидродинамическое давление, улучшая интегральные характеристики узла трения~\cite{gropper2016}.

Выделяют несколько основных механизмов влияния микрорельефа на характеристики смазочного слоя~\cite{fowell2007, gropper2016}.

Микрогидродинамический эффект является главным механизмом положительного влияния микрорельефа в условиях жидкостной смазки. Каждое углубление представляет собой миниатюрный клиновый зазор: на входной кромке (по ходу движения поверхности) зазор сужается, что приводит к локальному росту давления; на выходной кромке зазор расширяется, давление падает. Однако вследствие кавитации ($P \geq 0$) отрицательные давления обрезаются, и суммарный вклад каждого углубления в несущую способность оказывается положительным~\cite{fowell2007, brizmer2003}. Этот эффект наиболее выражен при неглубоких углублениях с отношением $h_p/c = 0{,}1\text{ - }0{,}3$~\cite{rahmani2007}.

Углубления выполняют функцию микрорезервуаров смазки. При пуске и останове машины, когда гидродинамический клин ещё не сформирован, смазка, удерживаемая в углублениях за счёт капиллярных сил, обеспечивает начальное смазывание контактных поверхностей и предотвращает сухое трение~\cite{galda2009}.

Углубления также служат ловушками для продуктов износа. Частицы изнашивания, попадающие в зазор, захватываются углублениями и выводятся из зоны контакта, что снижает абразивное воздействие на рабочие поверхности~\cite{etsion2005}.

Наиболее распространённой технологией создания микрорельефа является лазерное текстурирование поверхностей (Laser Surface Texturing, LST)~\cite{etsion2005}. Метод основан на локальном воздействии импульсного лазерного излучения, которое испаряет материал поверхности с высокой точностью (погрешность глубины менее 1~мкм). Помимо лазерного текстурирования применяются фрезерование на станках с ЧПУ, электроэрозионная обработка и накатывание~\cite{gropper2016}. Лазерный метод является наиболее перспективным благодаря высокой точности, воспроизводимости и возможности обработки локальных зон.

Углубления классифицируются по форме профиля: эллипсоидальные, цилиндрические, конические, сферические, параболические и комбинированные~\cite{gropper2016}. Эффективность текстурирования определяется совокупностью геометрических параметров: формой профиля, относительной глубиной $h_p/c$, размерами в плане ($a$, $b$), плотностью размещения и расположением текстурированной зоны относительно нагруженной области~\cite{tala2012, brizmer2003, rahmani2007}. Исследования~\cite{tala2012} показали, что текстурирование нагруженной зоны $90^\circ\text{ - }270^\circ$ даёт наибольший эффект по сравнению с полным текстурированием окружности.

Экспериментальные и численные исследования~\cite{kligerman2005, zhang2016, wang2013} подтвердили, что правильно подобранный микрорельеф способен снизить коэффициент трения на 10-30\% и увеличить несущую способность на 5-20\% в зависимости от условий работы. При этом наибольший положительный эффект наблюдается при высоких эксцентриситетах ($\varepsilon > 0{,}5$), когда минимальный зазор мал и относительный вклад углублений в формирование давления максимален~\cite{dobrica2010}.

\subsection*{Выводы по главе 1}
\phantomsection
\addcontentsline{toc}{subsection}{Выводы по главе 1}

Рассмотрены теоретические основы гидродинамической смазки подшипников скольжения. Описаны механизмы формирования несущего смазочного слоя, режимы смазки и основные интегральные характеристики. Проведён обзор исследований в области дискретного микрорельефа: описаны механизмы влияния микрорельефа на характеристики подшипника, технологии нанесения и классификация углублений по форме. Анализ литературы показал, что текстурирование рабочих поверхностей является эффективным способом повышения нагрузочной способности и снижения коэффициента трения.
